\documentclass[10pt,a4paper]{amsart}
\usepackage[margin=19mm]{geometry}
\usepackage{amsmath,amssymb,mathtools}
\usepackage[scr=rsfs,cal=euler]{mathalfa}
\usepackage{xfrac}
\usepackage[unicode,hidelinks,bookmarks=false]{hyperref}
\newcommand{\ie}{i.e.\ }
\newcommand{\cf}{cf.\ }
\newcommand{\resp}{respectively\ }
\newcommand{\Subsection}{\subsection}
\providecommand{\nobreakdash}{\nobreak\hbox{-}}
\setlength{\emergencystretch}{2em}
\multlinegap0pt
\usepackage{bm}
\usepackage{bbm}
\usepackage{dsfont}
\usepackage{xspace}
\usepackage{cancel}
\usepackage{mathtools}
\usepackage{amsthm}
\makeatletter
\@ifundefined{theorem}{\newtheorem{theorem}{Theorem}}{}
\@ifundefined{lemma}{\newtheorem{lemma}[theorem]{Lemma}}{}
\makeatother
\title[Free Boundary Plateau Model Cones in $\mathbb{B}^n$ are Rigi]{Free Boundary Plateau Model Cones in $\mathbb{B}^n$ are Rigid under Conformal Minimal Immersions}
\author{Elham Matinpour}
\date{Source: arXiv:2605.27776v1 (CC BY 4.0)}
\begin{document}
\maketitle

\noindent\emph{Source and licence.} Free Boundary Plateau Model Cones in $\mathbb{B}^n$ are Rigid under Conformal Minimal Immersions. Version arXiv:2605.27776v1 (CC BY 4.0), distributed under the Creative Commons Attribution 4.0 International licence, as recorded in the arXiv licence metadata for this exact version and shown on the abstract page https://arxiv.org/abs/2605.27776v1. Full rights notice: licenses/arxiv-2605.27776-CC-BY-4.0.txt.\par\smallskip

\noindent\textit{Editorial scope.} Counted unit: the Theorem whose source label is \texttt{None} in the article (source0.tex lines 874--881, source label \texttt{None}), delivered together with the article's own complete original proof of it (source0.tex lines 883--920, one \texttt{proof} environment of 38 lines). No mathematics is written, repaired or re-derived: statement and proof are the article's own text line for line. Required context retained uncounted: lem: uniq. of RT-network (lines 332--337). Every definition, display and label of those lines is retained unchanged. Omitted from this package and not redistributed: 1--331, 338--873, 882--882, 921--1110. Nothing inside a retained excerpt is omitted or altered: every retained body is delivered exactly as the article's lines are written, so no omission receipt of any kind accompanies this package. Citations: the retained excerpts carry no citation command at all, so no bibliography entry of the article is transcribed and no reference is redistributed. Reference adaptation: none was needed and none was made; every reference inside the delivered unit resolves inside it, so no unresolved reference or citation is produced. Printed slips kept literal: no printed word, symbol, hypothesis or number of the counted statement or of its proof was altered. Layout and class: the article's own class is \texttt{amsart} with packages including epsfig, color, blindtext, graphicx, enumitem, url, amssymb, import, comment, esint, xypic, amsthm, geometry, hyperref; the wrapper is the lane's pinned \texttt{amsart} editorial wrapper at 10pt on A4, which restates the theorem-like environments the retained text needs and adds the article's own macro definitions that the retained text needs (none), with extra wrapper packages (bm, bbm, dsfont, xspace, cancel, mathtools). The delivered numbering is the unit's own. Assets: no retained span contains a figure environment, a graphics-inclusion command, a PDF-page inclusion, a table or a TikZ picture; no image file is copied into this package, the package declares no asset dependency and no image file is redistributed with it. Licence: version arXiv:2605.27776v1 of this article is distributed under the Creative Commons Attribution 4.0 International licence, as recorded in the arXiv licence metadata for this exact version and shown on the abstract page https://arxiv.org/abs/2605.27776v1. A full rights notice is delivered at licenses/arxiv-2605.27776-CC-BY-4.0.txt. Characters: every retained excerpt is pure ASCII, so the delivered engine, XeLaTeX with its default Latin Modern text font, reports no missing character.
\begin{lemma}\label{lem: uniq. of RT-network}
Let $\mathcal{N} \subset \mathbb{S}^{n-1}$ be a stationary geodesic network with tetrahedral combinatorics, i.e., it has exactly four vertices, six geodesic edges, and every pair of vertices is connected by an edge. Then $\mathcal{N}$ is isometric to the regular tetrahedral network, that is, there exists an orthogonal transformation $A \in \mathrm{O}(n)$ such that
$
\mathcal{N} = A \bigl( \mathcal{NT} \bigr)$.\\
In particular, all edges of $\mathcal{N}$ have the same spherical length $\arccos(-1/3)$.
\end{lemma}

\begin{theorem}
Let \(\Sigma \subset B^n\) be a free-boundary minimal $2$-dimensional T-surface with exactly one T-point, four Y-junction edges (along each of which three faces meet at $120^\circ$), and six faces, each diffeomorphic to a quarter-disk. Assume that
\begin{itemize}
\item \(\partial \Sigma \cap \mathbb{S}^{n-1} \) is a stationary geodesic network with tetrahedral combinatorics,
\item \(\Sigma\) meets \(\mathbb{S}^{n-1}\) orthogonally.
\end{itemize}
Then there exists \(A \in \mathrm{O}(n)\) such that \(\Sigma = A(TC)\), where \(TC\) is the T-cone over the regular tetrahedron centered at the origin. In particular, the T-point is at the origin, the four Y-junctions are the straight radial segments from the origin to the vertices, and each face is a flat conical sector in the $2$-plane spanned by the corresponding pair of vertices.
\end{theorem}

\begin{proof}
Let \(N = \partial \Sigma \cap \mathbb{S}^{n-1}\) be the boundary network. By assumption \(N\) has tetrahedral combinatorics, each edge is a geodesic arc, and at each of its four vertices exactly three edges meet with the \(120^\circ\) equilibrium condition (the Y-configuration). Thus \(N\) is a stationary geodesic network with tetrahedral combinatorics. By the uniqueness lemma \ref{lem: uniq. of RT-network}, there exists \(A \in \mathrm{O}(n)\) such that \(N = A(\mathcal{RT})\), where \(\mathcal{RT}\) is the regular tetrahedral network on the great \(2\)-sphere in the subspace \(W \cong \mathbb{R}^3\) spanned by the standard regular-tetrahedron vertices \(v_1,\dots,v_4\) (with \(\langle v_i,v_j\rangle = -1/3\)). Without loss of generality (replacing \(\Sigma\) by \(A^{-1}(\Sigma)\)) we may assume that the boundary network is exactly the regular tetrahedral network \(\mathcal{RT} \subset W \cap \mathbb{S}^{n-1}\).
\medskip


Suppose $F$ is a face of $\Sigma$ with $\partial F=\sigma \cup \Gamma_1 \cup \Gamma_2$, where $\sigma$ is a geodesic arc on $\mathbb{RT}$ and $\Gamma_1,\Gamma_2$ are two junction curves of $\Sigma$. Let $\sigma\subset P$ be the plane passing through the origin. Then reflect $F$ and $P$ across the sphere along the boundary $\sigma$ via the inversion
$$
I(x)=\frac{x}{\vert x\vert^2}.
$$
Define $\widetilde{F}:=F\cup I(F).$
Then $\widetilde{F}$ is a smooth minimal surface in a neighborhood of $\sigma$, and $\sigma$ becomes an interior curve of $\widetilde{F}$. Since $P$ passes through the origin, we have $I(P)=P$. Thus, $P$ is also invariant under the reflection. By assumption, $\sigma \subset F \cap P$,
the outward conormal of $F$ agrees with that of $P$ along $\sigma$. This implies:
$$
T_pF=T_pP,\ \ \ \text{for all}\ p\in \sigma
$$
Thus, along $\sigma$, the surfaces $\widetilde{F}$ and $P$ intersect, and have the same tangent plane. So, they have second-order contact along $\sigma$. Therefore, by the unique continuation-strong maximum principle $\widetilde{F}\subset P$, and thus $F\subset P$.
\medskip

We know that there exists a regular $T$-cone $TC$ with $\partial TC=\mathbb{RT}$, and moreover, $TC$ has a flat face which passes through the origin and its boundary is $\sigma$. Since the argument holds for each face of $\Sigma$, and $\Sigma$ has the structure of a $T$-surface with a single $T$-point, four $Y$-junctions, and six faces, we conclude that $\Sigma=TC$. 
\medskip


Let \(p \in B^n\) be the unique T-point of \(\Sigma\). Near \(p\), \(\Sigma\) consists of six faces meeting along four Y-junction edges emanating from \(p\). Because \(\Sigma\) is minimal and the Y-junctions satisfy the \(120^\circ\) condition, the tangent cone of \(\Sigma\) at \(p\) is a stationary singular cone whose link (on the unit sphere in \(T_p\mathbb{R}^n\)) is a stationary geodesic network with tetrahedral combinatorics. Applying the uniqueness lemma once more to this link, the four outgoing unit tangent vectors \(\tau_1,\dots,\tau_4\) at \(p\) themselves form a regular tetrahedral configuration in some \(3\)-dimensional subspace \(V_p \subset T_p\mathbb{R}^n\), and the six local sectors are flat conical sectors in the six corresponding \(2\)-planes.
\medskip

The free-boundary condition (\(\Sigma\) meets \(\mathbb{S}^{n-1}\) orthogonally) together with the stationarity of the whole configuration implies that \(\Sigma\) cannot leave the \(3\)-dimensional subspace \(W\) containing the boundary: any component of a face transverse to \(W\) would produce a non-zero mean-curvature term or violate the balancing of conormals at the Y-junctions (by the maximum principle for minimal surfaces and the fact that the boundary data lie in \(W\)). Thus \(p \in W\) and \(V_p = W\) (the only \(3\)-dimensional subspace compatible with the boundary). Consequently the four outgoing directions \(\tau_k\) at \(p\) point precisely toward the four boundary vertices \(v_k\).
\medskip

Each Y-junction edge \(\gamma_k\) is therefore a curve from \(p\) to \(v_k\) whose initial direction is exactly the radial direction \(v_k - p\). By the triple-junction stationarity condition along \(\gamma_k\) (three minimal faces meeting at \(120^\circ\)) and the global minimality of \(\Sigma\), Lawlor's theorem (extended to free-boundary settings) forces each \(\gamma_k\) to be a straight line segment. Thus each \(\gamma_k\) is the radial segment \([p, v_k]\).
\medskip

It remains to show \(p = 0\). Suppose \(p \neq 0\). The free-boundary condition implies that the position vector field \(X(x) = x\) produces a vanishing first variation of area (the conormal along each boundary arc is parallel to the radial vector on \(\mathbb{S}^{n-1}\), and the junction balancing cancels the interior contributions). However, the monotonicity formula for free-boundary minimal surfaces in the ball shows that the area ratio \(r^{-2} \mathrm{Area}(\Sigma \cap B_r(0))\) is strictly increasing unless \(\Sigma\) is conical with vertex at \(0\). Since \(\Sigma\) has only one singularity (at \(p \neq 0\)) and its tangent cone at \(p\) is already conical, the only way the monotonicity ratio can be constant (forced by the first-variation vanishing for the radial field) is if \(p = 0\).
\medskip

With \(p = 0\) and the Y-junctions being the radial segments to the \(v_k\), each face is a minimal surface bounded by two radial segments and one great-circle arc on \(\mathbb{S}^{n-1}\). Because the face meets the sphere orthogonally and the two radial boundaries satisfy the \(120^\circ\) condition with the adjacent faces, the same holomorphic-quadratic-differential argument used for the Y-cone in Matinpour (2025) (applied sector by sector after conformal parametrization) shows that each face must be flat: the second fundamental form vanishes identically. Thus each face is precisely the flat conical sector in the plane spanned by the corresponding pair \(v_i,v_j\).

This is exactly the T-cone \(\mathcal{T}\). The proof is complete.
\end{proof}

\end{document}
